\section{Operators and Expressions in the Go Programming Language}

Go, commonly referred to as Golang, is a statically typed,
compiled programming language designed to support the development
of reliable, efficient and maintainable software. Operators and
expressions constitute fundamental elements of the language's
syntax and semantics. Operators specify operations performed on
operands, whilst expressions combine operands, operators and
function calls to produce values.

Go defines operators for arithmetic, comparison, logical
evaluation, bit manipulation, assignment, pointer operations and
other language-specific tasks. The validity of an expression
depends on the types of its operands, the operator being used
and the rules defined by the language specification.

\section{Operators in Go}

An operator is a syntactic element that specifies an operation
on one or more operands. Depending on the operator, the number
of operands may be one, two or, in certain syntactic forms,
associated with a specialised operation.

Operators in Go can be classified into several principal
categories.

\subsection{Arithmetic Operators}

Arithmetic operators perform numerical calculations. Go supports
the following operators:

\begin{center}
\begin{tabular}{|c|l|}
\hline
\textbf{Operator} & \textbf{Description} \\
\hline
\texttt{+} & Addition \\
\texttt{-} & Subtraction or unary negation \\
\texttt{*} & Multiplication \\
\texttt{/} & Division \\
\texttt{\%} & Remainder \\
\hline
\end{tabular}
\end{center}

The arithmetic operators \texttt{+}, \texttt{-}, \texttt{*}
and \texttt{/} are applicable to compatible numeric types.
The remainder operator \texttt{\%} is applicable to integer
operands.

For example:

\begin{lstlisting}
package main

import "fmt"

func main() {
    a := 17
    b := 5

    fmt.Println(a + b) // 22
    fmt.Println(a - b) // 12
    fmt.Println(a * b) // 85
    fmt.Println(a / b) // 3
    fmt.Println(a % b) // 2
}
\end{lstlisting}

When both operands of division are integers, the result is an
integer, with any fractional part discarded towards zero.
Consequently, the expression \texttt{17 / 5} evaluates to
\texttt{3}, rather than \texttt{3.4}.

Division by zero causes a run-time panic for integer operands.
Floating-point division follows IEEE 754 behaviour, allowing
results such as positive infinity or NaN in appropriate cases.

The addition operator also supports string concatenation.
For example, \texttt{"Go" + "Lang"} produces
\texttt{"GoLang"}.

\subsection{Comparison Operators}

Comparison operators compare two operands and produce a Boolean
result of type \texttt{bool}. The supported operators are:

\begin{center}
\begin{tabular}{|c|l|}
\hline
\textbf{Operator} & \textbf{Meaning} \\
\hline
\texttt{==} & Equal to \\
\texttt{!=} & Not equal to \\
\texttt{<} & Less than \\
\texttt{<=} & Less than or equal to \\
\texttt{>} & Greater than \\
\texttt{>=} & Greater than or equal to \\
\hline
\end{tabular}
\end{center}

The applicability of a comparison operator depends on the
operand types. Equality can be tested for comparable types,
including numeric values, strings, Booleans and pointers,
subject to the language's comparability rules. Ordered
comparisons are available for supported ordered types, such
as numbers and strings.

For example:

\begin{lstlisting}
x := 12
y := 20

fmt.Println(x < y)  // true
fmt.Println(x == y) // false
fmt.Println(x != y) // true
\end{lstlisting}

Go does not permit implicit conversion between different
numeric types. Therefore, comparing an \texttt{int} value
directly with a \texttt{float64} value generally requires an
explicit conversion.

\subsection{Logical Operators}

Logical operators operate on Boolean operands and are commonly
used to construct compound conditions.

\begin{center}
\begin{tabular}{|c|l|}
\hline
\textbf{Operator} & \textbf{Meaning} \\
\hline
\texttt{\&\&} & Logical AND \\
\texttt{||} & Logical OR \\
\texttt{!} & Logical NOT \\
\hline
\end{tabular}
\end{center}

The \texttt{\&\&} operator produces \texttt{true} only when
both operands are true. The \texttt{||} operator produces
\texttt{true} when at least one operand is true, whilst
\texttt{!} reverses a Boolean value.

For example:

\begin{lstlisting}
age := 25
hasID := true

eligible := age >= 18 && hasID
fmt.Println(eligible) // true
\end{lstlisting}

The logical operators \texttt{\&\&} and \texttt{||} use
short-circuit evaluation. The second operand is evaluated only
when necessary to determine the result. For instance, in an
expression of the form \texttt{A \&\& B}, operand \texttt{B}
is not evaluated if \texttt{A} is false.

\subsection{Bitwise and Bit-Shift Operators}

Bitwise operators manipulate individual bits of integer values.
They are useful in low-level programming, binary protocols,
bit masks and performance-sensitive operations.

\begin{center}
\begin{tabular}{|c|l|}
\hline
\textbf{Operator} & \textbf{Description} \\
\hline
\texttt{\&} & Bitwise AND \\
\texttt{|} & Bitwise OR \\
\texttt{\^{}} & Bitwise XOR, or unary bitwise complement \\
\texttt{\&\^} & Bit clear \\
\texttt{<<} & Left shift \\
\texttt{>>} & Right shift \\
\hline
\end{tabular}
\end{center}

The bit-clear operator \texttt{\&\^} clears those bits in its
left operand that are set in its right operand. The binary
operator \texttt{\^{}} performs exclusive OR, whereas its
unary form complements the bits of an integer operand.

For example:

\begin{lstlisting}
a := 12 // Binary: 1100
b := 10 // Binary: 1010

fmt.Println(a & b)   // 8
fmt.Println(a | b)   // 14
fmt.Println(a ^ b)   // 6
fmt.Println(a &^ b)  // 4
fmt.Println(a << 1)  // 24
fmt.Println(a >> 2)  // 3
\end{lstlisting}

Bitwise binary operators require integer operands of compatible
types. Shift operations move the bits of the left operand by
the number of positions specified by the right operand. The
shift count must be an integer expression, and a negative
shift count is not permitted.

\subsection{Assignment Operators}

Assignment operators associate values with variables or other
assignable operands. The simple assignment operator is
\texttt{=}, whilst compound assignment operators combine an
operation with assignment.

\begin{center}
\begin{tabular}{|c|l|}
\hline
\textbf{Operator} & \textbf{Equivalent operation} \\
\hline
\texttt{=} & Assignment \\
\texttt{+=} & Addition followed by assignment \\
\texttt{-=} & Subtraction followed by assignment \\
\texttt{*=} & Multiplication followed by assignment \\
\texttt{/=} & Division followed by assignment \\
\texttt{\%=} & Remainder followed by assignment \\
\texttt{\&=} & Bitwise AND followed by assignment \\
\texttt{|=} & Bitwise OR followed by assignment \\
\texttt{\^{ }=} & Bitwise XOR followed by assignment \\
\texttt{\&\^=} & Bit clear followed by assignment \\
\texttt{<<=} & Left shift followed by assignment \\
\texttt{>>=} & Right shift followed by assignment \\
\hline
\end{tabular}
\end{center}

For example:

\begin{lstlisting}
count := 10
count += 5
count *= 2

fmt.Println(count) // 30
\end{lstlisting}

The short declaration operator \texttt{:=} is distinct from
ordinary assignment. It declares one or more local variables
and initialises them, subject to the scope and declaration
rules of Go. It cannot be used as a general replacement for
\texttt{=} when all variables have already been declared.

\subsection{Pointer Operators}

Go provides two principal pointer operators:

\begin{itemize}
    \item \texttt{\&} obtains the address of an addressable
    operand.
    \item \texttt{*} dereferences a pointer to access the
    value stored at the referenced address.
\end{itemize}

For example:

\begin{lstlisting}
x := 42
p := &x

fmt.Println(*p) // 42

*p = 100
fmt.Println(x)  // 100
\end{lstlisting}

In this example, \texttt{p} stores a pointer to \texttt{x}.
Dereferencing \texttt{p} permits the referenced value to be
read or modified. Dereferencing a nil pointer causes a
run-time panic.

\subsection{Increment and Decrement}

Go supports the increment operator \texttt{++} and decrement
operator \texttt{--}. These statements increase or decrease
an integer or other suitable numeric variable by one.

\begin{lstlisting}
n := 5
n++
fmt.Println(n) // 6

n--
fmt.Println(n) // 5
\end{lstlisting}

Unlike languages such as C and C++, Go treats \texttt{++} and
\texttt{--} as statements rather than expressions. They cannot
be used as parts of larger expressions, and Go does not
support postfix or prefix value-producing forms such as
\texttt{++n} or \texttt{n++} within an expression.

\section{Expressions in Go}

An expression is a syntactic construct that represents a value
or computation. Expressions may consist of constants,
variables, operators, function calls, conversions and other
expressions combined according to the language grammar.

The type and value of an expression are determined by the
types of its components and the rules governing the operations
performed.

\subsection{Constant and Variable Expressions}

A constant expression represents a compile-time constant
value, whereas a variable expression refers to a variable
whose value may change during execution.

For example:

\begin{lstlisting}
const rate = 0.15
price := 200.0

discount := price * rate
finalPrice := price - discount
\end{lstlisting}

Here, \texttt{rate} is a constant and \texttt{price} is a
variable. The expressions calculate a discount and the
resulting price.

Go constants may be untyped until a context requires a
specific type. This permits certain constant expressions
to be evaluated with greater precision than ordinary
run-time floating-point expressions.

\subsection{Arithmetic Expressions}

Arithmetic expressions combine numeric operands with
arithmetic operators. Their validity depends on the
compatibility of operand types and the requirements of
the relevant operation.

For example:

\begin{lstlisting}
a := 8
b := 3
result := a*a + b*b
fmt.Println(result) // 73
\end{lstlisting}

The expression \texttt{a*a + b*b} consists of two
multiplication subexpressions combined by addition.
Operator precedence determines the order in which these
operations are grouped.

\subsection{Boolean Expressions}

Boolean expressions evaluate to values of type \texttt{bool}.
They commonly contain comparison and logical operators and
are used in conditional statements and loops.

\begin{lstlisting}
temperature := 22
isDry := true

if temperature >= 20 && isDry {
    fmt.Println("Suitable conditions")
}
\end{lstlisting}

The condition is true only when the temperature is at least
20 and \texttt{isDry} is true. Go requires conditions in
constructs such as \texttt{if} and \texttt{for} to be Boolean;
it does not implicitly convert integers or other values to
Boolean values.

\subsection{Function-Call Expressions}

A function-call expression invokes a function and may produce
one or more results. Function calls can appear within larger
expressions when their result types are appropriate.

\begin{lstlisting}
func square(n int) int {
    return n * n
}

func main() {
    result := square(6) + square(2)
    println(result) // 40
}
\end{lstlisting}

The expression \texttt{square(6) + square(2)} invokes the
function twice and adds the returned integer values.
A function returning multiple values generally requires a
context that accepts those results, such as a multiple
assignment or a return statement.

\subsection{Composite Expressions}

Composite expressions combine several operations, potentially
including arithmetic, comparisons, logical operators and
function calls. Parentheses can be used to make the intended
grouping explicit.

For example:

\begin{lstlisting}
a := 10
b := 4
c := 2

result := (a + b) * c
fmt.Println(result) // 28
\end{lstlisting}

The parentheses require the addition to be evaluated as
a subexpression before multiplication. Without parentheses,
operator precedence would group multiplication before
addition.

\section{Operator Precedence and Associativity}

Operator precedence determines how operators are grouped
when an expression contains multiple operators without
explicit parentheses. Go defines five precedence levels
for binary operators, listed below from highest to lowest.

\begin{center}
\begin{tabular}{|c|l|}
\hline
\textbf{Level} & \textbf{Operators} \\
\hline
5 & \texttt{* / \% << >> \& \&\^} \\
4 & \texttt{+ - | \^{}} \\
3 & \texttt{== != < <= > >=} \\
2 & \texttt{\&\&} \\
1 & \texttt{||} \\
\hline
\end{tabular}
\end{center}

Operators at higher precedence levels are grouped before
operators at lower levels. Operators within the same binary
precedence level associate from left to right.

For example:

\begin{lstlisting}
result := 2 + 3*4
fmt.Println(result) // 14
\end{lstlisting}

Multiplication has a higher precedence than addition;
therefore, the expression is grouped as
\texttt{2 + (3*4)}.

Parentheses should be used when they improve clarity,
particularly in expressions containing several operators.

\section{Type Rules and Conversions}

Go is statically typed. Consequently, the compiler checks
whether an expression uses operators with operands of
appropriate types. Unlike some dynamically typed languages,
Go does not generally perform implicit conversions between
different numeric types.

For example:

\begin{lstlisting}
var x int = 10
var y float64 = 2.5

// Invalid: mismatched operand types
// result := x + y

result := float64(x) + y
fmt.Println(result) // 12.5
\end{lstlisting}

The explicit conversion \texttt{float64(x)} converts the
integer value to \texttt{float64}, allowing the addition.
Conversions should be used with an understanding of their
effects, as converting floating-point values to integers
discards the fractional part, and converting between numeric
types can lose information.

\section{Evaluation and Run-time Behaviour}

Expression evaluation follows the language's specification
for operand evaluation, operator semantics and control flow.
Logical short-circuiting can prevent an operand from being
evaluated, whilst other expressions may invoke functions
that modify program state.

Programmers should not assume that arbitrary operands are
evaluated in an intuitively chosen order. Go specifies
particular evaluation-order rules, including left-to-right
ordering for function calls, method calls and communication
operations within an expression, as well as rules for
assignments and other constructs.

Certain expressions may lead to run-time panics, including
integer division by zero, nil-pointer dereferencing and
out-of-range indexing. These situations should be considered
when writing robust programs.

\section{Conclusion}

Operators and expressions are essential to Go's computational
model. Operators specify arithmetic, comparison, logical,
bitwise, assignment and pointer-related operations, whilst
expressions combine these operations to produce values and
control computational decisions.

A sound understanding of operator precedence, associativity,
type compatibility, explicit conversions and evaluation
behaviour is necessary for writing correct and maintainable
Go programs. Adhering to these rules improves code clarity
and reduces the likelihood of compilation errors and
unexpected run-time behaviour.

\section{Reference}

The authoritative source for the syntax and semantics of
Go operators and expressions is the official Go Language
Specification, particularly the sections on operators,
expressions, assignments, conversions and evaluation order.

\begin{itemize}
    \item \url{https://go.dev/ref/spec}
\end{itemize}

